%\documentclass[12pt,a4paper]{article}  % Use this line if this document will be released
\documentclass[12pt,a4paper,draft]{article}  % Use this line if this document is a draft
\usepackage{ifdraft}

%% Bibliography
\usepackage{etoolbox}
\newcommand{\bibfile}{\jobname.bib}
\newcommand{\universalbib}{ref.bib}
\ifdraft{\IfFileExists{\universalbib}{\renewcommand{\bibfile}{\universalbib}}{}}{}
\newcounter{cite}
\pretocmd{\cite}{\stepcounter{cite}}{}{}

%% Add line numbers in draft mode
\RequirePackage[mathlines]{lineno}
\ifdraft{\linenumbers}{}
\renewcommand{\linenumberfont}{\normalfont\scriptsize\sffamily\color{gray}}
\setlength{\linenumbersep}{\marginparsep}

%% Geometry
\usepackage[a4paper, textwidth=16.0cm, textheight=22.0cm]{geometry}
\renewcommand{\baselinestretch}{1.2}

%% Basic packages
\usepackage{amsmath,amsthm,amssymb,amsfonts}
\usepackage{mathtools}
\usepackage{empheq}
\usepackage{xcolor}
\usepackage[bbgreekl]{mathbbol}
\DeclareSymbolFontAlphabet{\mathbbm}{bbold}
\DeclareSymbolFontAlphabet{\mathbb}{AMSb}
\usepackage{bbm}
\usepackage{upgreek}
\usepackage{accents}
\usepackage{xspace}
\usepackage{rotating}
\usepackage{multirow,booktabs}
\usepackage[en-US]{datetime2}

%% Format of the table of content
\usepackage[normalem]{ulem}
\usepackage[toc,page]{appendix}
\renewcommand{\appendixpagename}{\Large{Appendix}}
\renewcommand{\appendixname}{Appendix}
\renewcommand{\appendixtocname}{Appendix}
\setcounter{tocdepth}{2}

%% Section title style
\usepackage{sectsty}
\sectionfont{\large}
\subsectionfont{\large}

%% Some colors
\definecolor{darkblue}{rgb}{0,0.1,0.5}
\definecolor{darkgreen}{rgb}{0,0.5,0.1}
\definecolor{darkyellow}{rgb}{0.65,0.65,0.01}

%% Todo notes
\ifdraft{
    \setlength{\marginparwidth}{2.42cm}
    \usepackage[tickmarkheight=3pt,textsize=small,backgroundcolor=blue!16,linecolor=purple,bordercolor=purple]{todonotes}
}{
    \newcommand{\todo}[1]{}
    \newcommand{\listoftodos}{}
}

%% Graph, tikz and pgf
\setlength{\unitlength}{1mm}
\usepackage{graphicx}
\usepackage{tikz,tikzscale,pgf}
\usetikzlibrary{arrows,arrows.meta,patterns,positioning,decorations.markings,shapes}
\usepackage{pgfplots}
\usepackage{pgfplotstable}
\usepackage[justification=centering]{caption}
\usepgfplotslibrary{fillbetween}
\pgfplotsset{compat=1.11}

%% Turn off some unharmful warnings in draft mode
\ifdraft{
    \usepackage{silence}
    \WarningFilter{xcolor}{Incompatible color definition on}
    \WarningFilter{hyperref}{Draft mode on}
    \WarningFilter{refcheck}{Unused label}
    \WarningFilter{microtype}{`draft' option active}
    \WarningFilter{latex}{Writing or overwriting file}
    \WarningFilter{latex}{Writing file}
    \WarningFilter{latex}{Tab has}
    \WarningFilter{latex}{Marginpar on page}
    \WarningFilter{latex}{author given}
}{}

%% Hyperref, url, and email
\ifdraft{\usepackage{refcheck}\newcommand{\url}{\texttt}}{
    \usepackage{hyperref}
    \hypersetup{colorlinks, linkcolor=darkblue, anchorcolor=darkblue, citecolor=darkblue, urlcolor=darkblue}
    \usepackage{url}
}
\newcommand{\email}{\texttt}

%% Enumerate and itemize
\usepackage{eqlist}
\usepackage{enumitem}
\setlist[itemize]{leftmargin=*}
\setlist[enumerate]{leftmargin=*,label=\normalfont{(\alph*)}}

%% Algorithm environment
\usepackage[section]{algorithm}
\usepackage{algpseudocode,algorithmicx}
\newcommand{\INPUT}{\textbf{Input}}
\newcommand{\FOR}{\textbf{For}~}
\algrenewcommand\algorithmicrequire{\textbf{Input}}
\algrenewcommand\algorithmicensure{\textbf{Output}}
\algrenewcommand\alglinenumber[1]{\normalsize #1.}
\newcommand*\Let[2]{\State #1 $=$ #2}

%% Theorem-like environments
\newtheorem{theorem}{Theorem}[section]
\newtheorem{conjecture}{Conjecture}[section]
\newtheorem{corollary}{Corollary}[section]
\newtheorem{exercise}{Exercise}[section]
\newtheorem{lemma}{Lemma}[section]
\newtheorem{problem}{Problem}[section]
\newtheorem{proposition}{Proposition}[section]
\newtheorem{assumption}{Assumption}[section]
\newtheorem{example}{Example}[section]
\newtheorem{question}{Question}[section]
\theoremstyle{definition}
\newtheorem{definition}{Definition}[section]
\newtheorem{remark}{Remark}[section]
\usepackage{xpatch}
\xpatchcmd{\proof}{\itshape}{\normalfont\proofnamefont}{}{}
\newcommand{\proofnamefont}{\bfseries}

%% Equation numbering
\numberwithin{equation}{section}

%% Fine tuning
\usepackage{microtype}
\usepackage[nobottomtitles*]{titlesec}
\interfootnotelinepenalty=10000
\interdisplaylinepenalty=10000
\relpenalty=10000
\binoppenalty=10000
\clubpenalty=10000
\widowpenalty=10000
\displaywidowpenalty=10000

\mathcode`@="8000{\catcode`\@=\active\gdef@{\mkern1mu}}

%% Operators, commands
\usepackage{relsize}
\usepackage{nccmath}
\DeclareMathOperator*{\mcap}{\,\mathsmaller{\bigcap}\,}
\DeclareMathOperator*{\mcup}{\,\mathsmaller{\bigcup}\,}
\newcommand{\ceil}[1]{ {\lceil{#1}\rceil} }
\newcommand{\floor}[1]{ {\lfloor{#1}\rfloor} }

\DeclareMathOperator{\tr}{tr}
\DeclareMathOperator{\sort}{sort}
\DeclareMathOperator*{\Argmax}{Argmax}
\DeclareMathOperator*{\Argmin}{Argmin}
\DeclareMathOperator*{\Arglocmin}{Arglocmin}
\DeclareMathOperator*{\argmax}{argmax}
\DeclareMathOperator*{\argmin}{argmin}
\DeclareMathOperator*{\diag}{diag}
\DeclareMathOperator*{\Diag}{Diag}
\DeclareMathOperator{\Span}{span}
\DeclareMathOperator{\med}{med}
\DeclareMathOperator{\essinf}{essinf}
\DeclareMathOperator{\cl}{cl}
\DeclareMathOperator{\vol}{vol}
\DeclareMathOperator{\comp}{C}
\DeclareMathOperator{\sign}{sign}
\DeclareMathOperator{\rank}{rank}
\DeclareMathOperator{\range}{range}
\DeclareMathOperator{\card}{card}
\DeclareMathOperator{\diam}{diam}
\DeclareMathOperator{\dist}{dist}
\newcommand{\disth}{{\operatorname{\updelta_{\sss{H}}}}}
\newcommand{\ind}{\mathbbm{1}}
\newcommand\defeq{\mathrel{\overset{\makebox[0pt]{\mbox{\normalfont\tiny\sffamily def}}}{=}}}

\newcommand{\RR}{\mathbb{R}}
\newcommand{\BB}{\mathcal{B}}
\renewcommand{\SS}{\mathbb{S}}
\newcommand{\TT}{\mathcal{T}}
\newcommand{\ZZ}{\mathbb{Z}}
\newcommand{\NN}{\mathbb{N}}
\newcommand{\FF}{\mathcal{F}}
\newcommand{\CC}{\mathbb{C}}
\newcommand{\XX}{\mathcal{X}}
\newcommand{\sset}{\mathcal{S}}
\newcommand{\pen}{h}
\newcommand{\penpar}{\mu}
\newcommand{\res}{\rho}
\newcommand{\col}{r}
\newcommand{\ofd}{\mathcal{F}}
\newcommand{\stf}[1]{\mathbb{S}^{#1}}
\newcommand{\sss}[1]{{\scriptscriptstyle{#1}}}
\newcommand{\sK}{{\scriptscriptstyle{K}}}
\newcommand{\sT}{{\scriptscriptstyle{T}}}
\newcommand{\fro}{{\scriptstyle{\textnormal{F}}}}
\newcommand{\trs}{{\scriptstyle{\mathsf{T}}}}
\newcommand{\hmt}{{\scriptstyle{{\mathsf{H}}}}}
\newcommand{\pin}{{\scriptstyle{{\mathsf{+}}}}}
\newcommand{\inv}{{-1}}
\newcommand{\adj}{*}
\newcommand{\ones}{\mathbf{1}}

\newcommand{\cs}{\text{c}}
\newcommand{\hp}{\circ}
\newcommand{\cc}{\sss{\textnormal{C}}}
\newcommand{\dec}{\sss{\textnormal{D}}}
\newcommand{\cauchy}{\sss{\textnormal{C}}}
\newcommand{\scauchy}{\sss{\textnormal{S}}}
\newcommand{\crit}{\textnormal{crit}}
\newcommand{\rsg}{\hat{\partial}}
\newcommand{\gsg}{\partial}
\newcommand{\dom}{\textnormal{dom}}
\newcommand{\tf}{{\textnormal{f}}}
\newcommand{\tg}{{\textnormal{g}}}
\newcommand{\ts}{{\textnormal{s}}}
\newcommand{\st}{\textnormal{s.t.}}
\newcommand{\etc}{{etc.}\xspace}
\newcommand{\ie}{{i.e.}\xspace}
\newcommand{\eg}{{e.g.}\xspace}
\newcommand{\etal}{{et al.}\xspace}
\newcommand{\iid}{\text{i.i.d.}\xspace}
\newcommand{\as}{\text{a.s.}\xspace}

\newcommand{\me}{\mathrm{e}}
\newcommand{\md}{\mathrm{d}}
\newcommand{\mi}{\mathrm{i}}
\newcommand{\lev}{\mathrm{lev}}
\newcommand{\bA}{\mathbf{A}}
\newcommand{\bx}{\mathbf{u}}
\newcommand{\bb}{\mathbf{r}}
\newcommand{\nov}{n_{\textnormal{o}}}
\xspaceaddexceptions{]\}}
\newcommand{\MATLAB}{\textsc{Matlab}\xspace}
\newcommand{\octave}{\mbox{GNU Octave}\xspace}
\newcommand{\prblm}{\texttt}
\DeclareMathAlphabet{\mathsfit}{T1}{\sfdefault}{\mddefault}{\sldefault}
\SetMathAlphabet{\mathsfit}{bold}{T1}{\sfdefault}{\bfdefault}{\sldefault}
\newcommand{\prbb}{\mathsfit{p}}
\newcommand{\pp}{\mathsf{p}}
\newcommand{\qq}{\mathsf{q}}
\newcommand{\ttt}{\mathsfit{t}}
\newcommand{\tol}{\varepsilon}
\newcommand{\bt}{\mathbf{t}}
\newcommand{\br}{\mathbf{r}}
\newcommand{\dd}{\mathbf{d}}
\newcommand{\ii}{\mathbf{i}}
\newcommand{\jj}{\mathbf{j}}
\newcommand{\xx}{\mathbf{x}}
\renewcommand{\pp}{\mathbf{p}}
\renewcommand{\ggg}{\mathbf{g}}
\newcommand{\GG}{\mathbf{G}}
\DeclareMathOperator{\expc}{\mathbb{E}}
\renewcommand{\Pr}{\mathbb{P}}
\newcommand{\lb}{\underline}
\newcommand{\ub}{\overline}

\DeclareFontFamily{U}{dutchcal}{\skewchar\font=45 }
\DeclareFontShape{U}{dutchcal}{m}{n}{<-> s*[1.0] dutchcal-r}{}
\DeclareFontShape{U}{dutchcal}{b}{n}{<-> s*[1.0] dutchcal-b}{}
\DeclareMathAlphabet{\mathlcal}{U}{dutchcal}{m}{n}
\SetMathAlphabet{\mathlcal}{bold}{U}{dutchcal}{b}{n}
\usepackage[scr=boondoxupr]{mathalfa}
\newcommand{\model}{\tilde{f}}
\newcommand{\rmod}{F}
\newcommand{\Set}[1]{\mathcal{#1}}
\DeclareMathAlphabet{\mathpzc}{OT1}{pzc}{m}{it}
\newcommand{\slv}{\mathpzc}
\newcommand{\software}{\texttt}
\DeclareMathOperator{\eff}{\mathsf{e}\;\!}
\DeclareMathOperator{\Eff}{\mathsf{E}\;\!}
\newcommand{\out}{{\text{out}}}

%% Commands for revision
\newcommand{\red}[1]{\textcolor{red}{#1}}
\newcommand{\blue}[1]{\textcolor{blue}{#1}}
\newcommand{\green}[1]{\textcolor{darkgreen}{#1}}
\newcommand{\TYPO}[1]{{\color{orange}{#1}}}
\newcommand{\MISTAKE}[1]{{\color{violet}{#1}}}
\newcommand{\REPHRASE}[1]{{\color{darkgreen}{#1}}}
\newcommand{\REVISE}[1]{{\color{blue}{#1}}}
\newcommand{\REVISEred}[1]{{\color{red}{#1}}}
\newcommand{\COMMENT}{\todo}

%%%%%%%%%%%%%%%%%%%%%%%%%%%%%%%%%%%%%%%%%%%%%%%%%%%%%%%%%%%%%%%%%%%%%%%%%%%%%%%%%%%%%%%%%%%%%%%%%%%%
\title{Course Notes on Mathematical Analysis}
\date{\DTMnow}
\author{
    Zhu Huatao
    \thanks{}
}
%%%%%%%%%%%%%%%%%%%%%%%%%%%%%%%%%%%%%%%%%%%%%%%%%%%%%%%%%%%%%%%%%%%%%%%%%%%%%%%%%%%%%%%%%%%%%%%%%%%%

\begin{document}
\maketitle

\section{Real Number}

In this section we construct the real number system~$\RR$ from rational numbers.
Then we establish the completeness of~$\RR$ and present several classical theorems.
These theorems are equivalent forms of completeness and they are basic tools in analysis.

\subsection{Dedekind Cut}

We start from rational numbers.
They form an ordered field, and they are dense in themselves.
Still, they are not complete.
This gap motivates the construction of~$\RR$.

\begin{definition}[Rational Number]
A rational number is a number that can be written as a quotient of two integers.
The denominator is positive.
The numerator and denominator are coprime.
\[
\mathbb{Q}
=
\left\{
\frac{p}{q}
\;\middle|\;
p \in \ZZ,\ q \in \NN_{>0},\ (|p|,q)=1
\right\}.
\]
\end{definition}

\begin{definition}[Cut]
A cut is a subset~$A \subsetneq \mathbb{Q}$ with the following properties
\begin{enumerate}
\item~$A \neq \varnothing$
\item If~$x \in A$ and~$y < x$ with~$y \in \mathbb{Q}$, then~$y \in A$
\item~$A$ has no greatest element
\end{enumerate}
\end{definition}

A cut describes all rational numbers below one location on the number line.
If that location is rational then the cut corresponds to that rational number.
If that location is irrational then the cut still makes sense.
This is the key idea that allows us to define real numbers.

\begin{definition}[Order and Arithmetic on Cuts]
Let~$A$ and~$B$ be cuts.

\begin{enumerate}
\item We write~$A \le B$ if~$A \subseteq B$
\item We write~$A < B$ if~$A \subsetneq B$
\item We define addition by
\[
A+B
=
\{a+b \mid a\in A,\ b\in B\}
\]
\item If~$A \ge 0$ and~$B \ge 0$, we define multiplication by
\[
A\cdot B
=
\{ab \mid a\in A,\ b\in B,\ a>0,\ b>0\}\cup\mathbb{Q}_{\le0}
\]
\item For a cut~$A$ we define its additive inverse~$-A$ as the unique cut satisfying~$(-A)+A=0$
\end{enumerate}
\end{definition}

The next step is to check that these constructions are well defined.

\begin{proposition}
If~$A$ and~$B$ are cuts then~$A+B$ is a cut.
If~$A,B\ge0$ then~$A\cdot B$ is a cut.
\end{proposition}

\begin{proof}
We first treat~$A+B$.

It is non empty since~$A$ and~$B$ are non empty.
Choose~$a\in A$ and~$b\in B$.
Then~$a+b\in A+B$.
It is not all of~$\mathbb{Q}$ since~$A$ and~$B$ are proper lower sets.
There exist~$\alpha\notin A$ and~$\beta\notin B$.
If~$x\in A+B$, write~$x=a+b$ with~$a\in A$ and~$b\in B$.
Then~$a<\alpha$ and~$b<\beta$, hence~$x<\alpha+\beta$.
So~$\alpha+\beta\notin A+B$.

Now let~$x\in A+B$ and~$y<x$.
Take~$x=a+b$.
Set~$\delta=x-y>0$ and define~$a'=a-\delta$.
Then~$a'<a$, so~$a'\in A$ by downward closure.
Also~$y=a'+b$.
Hence~$y\in A+B$.

Finally~$A+B$ has no greatest element.
Take~$x=a+b\in A+B$.
Since~$A$ has no greatest element, choose~$a_1\in A$ with~$a<a_1$.
Then~$x<a_1+b\in A+B$.

So~$A+B$ is a cut.

For~$A\cdot B$ with~$A,B\ge0$ the argument is parallel.
It is non empty because~$\mathbb{Q}_{\le0}\subseteq A\cdot B$.
It is proper because positive factors are bounded above by numbers not in~$A$ and~$B$.
It is downward closed by construction and positivity.
It has no greatest element since for positive part one enlarges one factor inside its cut.
Hence~$A\cdot B$ is a cut.
\end{proof}

\begin{proposition}
The set of all cuts with the operations above is an ordered field.
\end{proposition}

\begin{proof}
Order properties follow from inclusion.
Reflexivity, antisymmetry, and transitivity are immediate.
Totality follows from the cut structure since for two cuts~$A$ and~$B$, either~$A\subseteq B$ or~$B\subseteq A$.

For addition and multiplication one checks closure by the previous proposition.
Associativity and commutativity follow from corresponding properties in~$\mathbb{Q}$.
Neutral elements are the cuts corresponding to~$0$ and~$1$.
Additive inverses exist by definition.
Multiplicative inverses for positive nonzero cuts are constructed by reciprocal lower sets and then extended by sign.
Distributivity follows from rational distributivity and set definitions.
Hence this system is an ordered field.
\end{proof}

We now define real numbers.

\begin{definition}[Real Number]
The set of real numbers~$\RR$ is the set of all Dedekind cuts of~$\mathbb{Q}$.
\end{definition}

This construction explains why~$\RR$ extends~$\mathbb{Q}$ and fills every gap.
The remaining task is to state and prove completeness in usable forms.

\subsection{Equivalent Forms of Completeness}

We list six classical statements.
Later we show they are equivalent.

\begin{theorem}[Supremum Principle]\label{supremum}
Every non empty subset of~$\RR$ that is bounded above has a least upper bound in~$\RR$.
\end{theorem}

\begin{theorem}[Monotone Convergence Theorem]\label{monotone}
Every monotone sequence in~$\RR$ that is bounded converges in~$\RR$.
\end{theorem}

\begin{theorem}[Nested Interval Theorem]\label{nested}
Let~$I_n=[a_n,b_n]$ with~$I_{n+1}\subseteq I_n$ for all~$n$.
If~$\lim_{n\to\infty}(b_n-a_n)=0$, then there is exactly one~$x\in\RR$ with~$x\in I_n$ for all~$n$.
\end{theorem}

\begin{theorem}[Bolzano Weierstrass Theorem]\label{bolzano}
Every bounded sequence in~$\RR$ has a convergent subsequence.
\end{theorem}

\begin{theorem}[Cauchy Criterion]\label{cauchy}
A sequence in~$\RR$ converges if and only if it is a Cauchy sequence.
\end{theorem}

\begin{theorem}[Heine Borel Theorem]\label{heine-borel}
A subset of~$\RR$ is compact if and only if it is closed and bounded.
\end{theorem}

\begin{proof}
We first prove Theorem~\ref{supremum} from Dedekind cuts.
Then we prove a cycle of implications among Theorem~\ref{supremum} to Theorem~\ref{heine-borel}.

\textbf{Step 1  The Dedekind Cuts imply the Supremum Principle}

Let~$S\subseteq\RR$ be non empty and bounded above.
Define
\[
U
=
\bigcup_{x\in S}x
\]
where each~$x$ is a cut.

We show~$U$ is a cut.
It is non empty since~$S$ is non empty and every cut is non empty.
It is proper since~$S$ has an upper bound~$M$, so each~$x\in S$ satisfies~$x\subseteq M$, hence~$U\subseteq M\subsetneq\mathbb{Q}$.
It is downward closed because each~$x$ is downward closed.
It has no greatest element since if~$u\in U$, then~$u\in x_0$ for some~$x_0\in S$, and~$x_0$ has no greatest element.

So~$U\in\RR$.
Now every~$x\in S$ satisfies~$x\subseteq U$, so~$x\le U$.
Thus~$U$ is an upper bound.
If~$V$ is any upper bound of~$S$, then~$x\subseteq V$ for all~$x\in S$, so~$U\subseteq V$.
Hence~$U\le V$.
Therefore~$U=\sup S$.

\textbf{Step 2  The Supremum Principle implies the Monotone Convergence Theorem}

Let~$(x_n)$ be increasing and bounded above.
Set~$S=\{x_n\mid n\in\NN\}$.
By supremum principle let~$x=\sup S$.
Given~$\varepsilon>0$,~$x-\varepsilon$ is not an upper bound of~$S$, so there exists~$N$ with~$x_N>x-\varepsilon$.
For~$n\ge N$ we have~$x-\varepsilon<x_N\le x_n\le x<x+\varepsilon$.
Hence~$|x_n-x|<\varepsilon$ for~$n\ge N$.
So~$x_n\to x$.

\textbf{Step 3  The Monotone Convergence Theorem implies the Nested Interval Theorem}

Define~$a_n$ and~$b_n$ from~$I_n=[a_n,b_n]$.
Since~$I_{n+1}\subseteq I_n$, sequence~$(a_n)$ is increasing and~$(b_n)$ is decreasing.
Also~$a_n\le b_n$.
By monotone convergence there exist~$a=\lim a_n$ and~$b=\lim b_n$.
Given~$b_n-a_n\to0$, we get~$b-a=0$, so~$a=b$.
Call this value~$x$.
Then~$a_n\le x\le b_n$ for all~$n$, hence~$x\in I_n$ for all~$n$.
Uniqueness follows since interval diameters tend to~$0$.

\textbf{Step 4  The Nested Interval Theorem implies the Bolzano Weierstrass Theorem}

Let~$(x_n)$ be bounded.
Choose~$[A_1,B_1]$ containing all~$x_n$.
Bisect it.
At least one half contains infinitely many terms.
Call that half~$I_1$.
Repeat recursively and obtain~$I_k=[A_k,B_k]$ with
\[
I_{k+1}\subseteq I_k,\qquad B_k-A_k=\frac{B_1-A_1}{2^{k-1}}
\]
and each~$I_k$ contains infinitely many terms.
By nested interval theorem there exists unique~$x\in\bigcap_k I_k$.
Choose~$n_k$ strictly increasing with~$x_{n_k}\in I_k$.
Then~$|x_{n_k}-x|\le B_k-A_k\to0$.
So~$x_{n_k}\to x$.

\textbf{Step 5  The Bolzano Weierstrass Theorem implies the Cauchy Criterion}

If~$x_n\to x$, then~$(x_n)$ is Cauchy.
This direction is standard.

Now suppose~$(x_n)$ is Cauchy.
Every Cauchy sequence is bounded.
By Bolzano Weierstrass there is a subsequence~$x_{n_k}\to x$.
Given~$\varepsilon>0$, choose~$N_1$ so that~$m,n\ge N_1$ implies~$|x_m-x_n|<\varepsilon/2$.
Choose~$K$ so that~$k\ge K$ implies~$n_k\ge N_1$ and~$|x_{n_k}-x|<\varepsilon/2$.
For~$n\ge N_1$ and~$k\ge K$
\[
|x_n-x|
\le
|x_n-x_{n_k}|+|x_{n_k}-x|
<
\varepsilon
\]
Hence~$x_n\to x$.

\textbf{Step 6  The Cauchy Criterion implies the Heine Borel Theorem}

Assume every Cauchy sequence converges.

First show closed and bounded implies compact.
Let~$K\subseteq\RR$ be closed and bounded.
Take any sequence in~$K$.
Boundedness and Bolzano Weierstrass can be recovered from Cauchy criterion by interval bisection and Cauchy subsequences.
Hence there is a convergent subsequence with limit in~$\RR$.
Closedness gives limit in~$K$.
So~$K$ is sequentially compact.
In metric spaces sequential compactness and compactness are equivalent.
Thus~$K$ is compact.

Now show compact implies closed and bounded.
If~K is compact then it is bounded, otherwise one constructs a sequence escaping to infinity with no convergent subsequence.
It is closed since limits of convergent sequences in~K stay in~K by compactness and uniqueness of limits.
Hence~K is closed and bounded.

So~K is compact if and only if it is closed and bounded.

\textbf{Step 7  The Heine Borel Theorem implies the Supremum Principle}

Let~$S\subseteq\RR$ be non empty and bounded above.
Pick~$a\in S$ and an upper bound~$b$.
Set~$K=[a,b]$.
By Heine Borel,~K is compact.

Define
\[
U=\{x\in K\mid x\ \text{is an upper bound of }S\}
\]
Then~U is non empty and closed in~K.
Define
\[
L=\{x\in K\mid x\le s\ \text{for some }s\in S\}
\]
Then~L is non empty and closed in~K.
Also every point of~K belongs to~L or~U.
By order structure there is a boundary point~c between these two sets.
One checks~c is an upper bound of~S and every smaller number fails to be an upper bound.
Hence~c=$\sup S$.

This completes the cycle and proves equivalence of the six theorems.
\end{proof}

%%%%%%%%%%%%%%%%%%%%%%%%%%%%%%%%%%%%%%%%%%%%%%%%%%%%%%%%%%%%%%%%%%%%%%%%%%%%%%%%%%%%%%%%%%%%%%%%%%%%
%% References
\ifnum\value{cite}>0
    \small
    \bibliography{\bibfile}
    \bibliographystyle{plain}
\fi

\end{document}